\documentclass[10pt,a4paper]{amsart}
\usepackage{amsmath,amssymb,amsthm,graphicx}
\usepackage{amsmath}
\usepackage{amssymb}
\usepackage{hyperref}
\newtheorem{Theorem}{Theorem}[section]
\newtheorem{Corollary}[Theorem]{Corollary}
\newtheorem{Lemma}[Theorem]{Lemma}
\newtheorem{Proposition}[Theorem]{Proposition}
{ \theoremstyle{definition}
	\newtheorem{Definition}[Theorem]{Definition}
	\newtheorem{Note}[Theorem]{Note}
	\newtheorem{Example}[Theorem]{Example}
	\newtheorem{Remark}[Theorem]{Remark} }

\usepackage[all]{xy}


\def\KK{\ensuremath\mathbb{K}}
\def\MM{\ensuremath\mathbb{M}}
\def\BB{\ensuremath\mathbb{B}}
\def\GL{\ensuremath\mathbf{GL}}
\def\UU{\ensuremath\mathbf{U}}
\def\PP{\ensuremath\mathbb{P}}
\def\CC{\ensuremath{\mathbb{C}}}
\def\NN{\ensuremath{\mathbb{N}}}
\def\QQ{\ensuremath{\mathbb{Q}}}
\def\RR{\ensuremath{\mathbb{R}}}
\def\ZZ{\ensuremath{\mathbb{Z}}}

\newcommand{\profinQ}{\mathfrak{Q}}
\newcommand{\reals}{\mathbb{R}}
\newcommand{\complexs}{\mathbb{C}}
\newcommand{\naturals}{\mathbb{N}}
\newcommand{\integers}{\mathbb{Z}}
\newcommand{\rationals}{\mathbb{Q}}
\newcommand{\quaternions}{\mathbb{H}}
\newcommand{\hyperbolic}{\mathbb{H}}
\newcommand{\K}{\mathbb{K}}
\newcommand{\realsorcomplexs}{\mathbb{K}}
\newcommand{\finitefield}{\mathbb{F}}
\newcommand{\Proj}{\mathbb{P}}
\DeclareMathOperator{\id}{id}
\DeclareMathOperator{\Ind}{Ind}
\newcommand{\boundary}[1]{\partial#1}
\newcommand{\boundedops}{\mathcal{B}}
\newcommand{\bundleover}{\!\!\downarrow\!\!}
\newcommand{\abs}[1]{\left\lvert#1\right\rvert} %absolute value
%\newcommand{\bigabs}[1]{\left\lvert#1\right\rvert} %absolute value
\newcommand{\norm}[1]{\left\lVert#1\right\rVert}
\newcommand{\generate}[1]{\langle#1\rangle}
\newcommand{\tensor}{\otimes}
\newcommand{\into}{\hookrightarrow}
\newcommand{\onto}{\twoheadrightarrow}
\newcommand{\iso}{\cong}
\newcommand{\disjointunion}{\amalg}
\newcommand{\subgroup}{\leq}
\newcommand{\supergroup}{\geq}
\newcommand{\superset}{\supset}
\newcommand{\normalsubgroup}{\lhd}
\newcommand{\semiProd}{\rtimes}
\newcommand{\semiprod}{\semiProd}
\newcommand{\dividesnot}{\nmid}
\newcommand{\divides}{\mid}
\DeclareMathOperator{\Pos}{Pos}
\DeclareMathOperator{\Aut}{Aut}
\DeclareMathOperator{\supp}{supp} %support
\DeclareMathOperator{\clos}{clos} %closure
\DeclareMathOperator{\im}{im} %image
\DeclareMathOperator{\vol}{vol} %Volume
\DeclareMathOperator{\diam}{diam} %Diameter
\DeclareMathOperator{\dist}{dist} %Distance
\DeclareMathOperator{\ord}{ord} %order
\DeclareMathOperator{\End}{End} %Endomorphisms
\DeclareMathOperator{\Hom}{Hom} %Homomorphisms
\DeclareMathOperator{\Span}{span}
\DeclareMathOperator{\spin}{spin}
\DeclareMathOperator{\diag}{diag}
\DeclareMathOperator{\spec}{Spec}
\DeclareMathOperator{\trace}{trace}
\DeclareMathOperator{\tr}{tr}
\DeclareMathOperator{\pr}{pr}
%\DeclareMathOperator{\res}{res} %Restriction
\DeclareMathOperator{\grad}{grad} %Restriction
% \DeclareMathOperator{\Sp}{Sp} %Spur % schon definiert???
\DeclareMathOperator{\fl}{f\/l} %flip
\DeclareMathOperator{\coker}{coker}
\DeclareMathOperator{\Char}{Char}
\DeclareMathOperator{\ind}{ind}
\DeclareMathOperator{\sgn}{sgn}
\DeclareMathOperator{\Ric}{Ric} % Ricci curvature
\DeclareMathOperator{\scal}{scal} % Scalar curvature
\DeclareMathOperator*{\dirlim}{dirlim}
\DeclareMathOperator*{\invlim}{invlim}
\DeclareMathOperator{\rank}{rank}
\DeclareMathOperator{\coind}{coind}
\DeclareMathOperator{\res}{res}
\newcommand{\bigdisjointunion}{\bigsqcup}
\DeclareMathOperator{\Orb}{Orb}
\DeclareMathOperator{\GTOP}{G-TOP}
\newcommand{\forget}[1]{}
\newenvironment{Forget}{}{}
\newcommand{\innerprod}[1]{\langle #1 \rangle}
\title{The equivariant Stolz exact sequence descends to the quotient by the fundamental group}
\author{Massimiliano Puglisi, Thomas Schick, Vito Felice Zenobi}
\date{doi:10.3842/SIGMA.2025.093; CC BY 4.0}
\begin{document}
\maketitle
\noindent\textit{Editorial note.} All delivered mathematical text is the author's own, copied from the pinned publisher source sources/geometry-fl04-stolz-positive-scalar-curvature-st/pinned-original-source.tex; no statement or proof step was rewritten, reordered or removed.\par\medskip
\noindent\textit{Editorial note.} Preamble provenance: the retained preamble is the source's own layout adapter, if any, plus the author's macro block at lines 11-108 of sources/geometry-fl04-stolz-positive-scalar-curvature-st/pinned-original-source.tex (everything between the publisher class metadata and \textbackslash{}begin\{document\}, comments excluded). Nothing was added, no macro was invented, and the publisher document class is deliberately not reproduced; the excerpt serves only as preamble, never as a span coordinate.\par\medskip
\noindent\textit{Editorial note.} Metadata provenance: source\_id and source\_version were taken from sources/geometry-fl04-stolz-positive-scalar-curvature-st/license-observation.json; the title from SOURCE.json; the author list from the source metadata record; the source URL from the first non-empty of source\_url, source\_urls[0] and the evidence's fixed\_version\_url/license\_url. The licence evidence file used by this card is sources/geometry-fl04-stolz-positive-scalar-curvature-st/license-observation.json, recorded in license.evidence\_path so that the binding target is explicit.\par\medskip
\noindent\textit{Editorial note.} External references are navigation pointers only; every deliberately kept step of the author's argument stays inside the counted spans. The displayed numbers of those pointers are the author's own numbering under the author's \textbackslash{}newtheorem declarations (section.index as declared), computed mechanically by the lane's scanner.\par\medskip
\noindent\textit{Editorial note.} The counted proof contains no reference to any labelled item and no bibliography citation, so the card needs neither a navigation stub nor a bibliography excerpt.\par\medskip
\noindent\textit{Editorial note.} This is the first card this lane authors from this source; it was picked as a zero-stub candidate (no navigation pointer at all) following the agreed ordering.\par\medskip
\noindent\textit{Allowed standard prerequisites (editorial).}\begin{itemize}
\item The G-equivariant and G-tilde-equivariant Stolz exact sequences attached to X and its covering X-tilde. Set up in the paragraphs preceding the counted statement, ending at source line 509. Fixes the two objects whose isomorphism is asserted and the notion of equivariant cycle and bordism used in the proof.
\item The covering X-tilde of X with group of deck transformations pi\_1(X) inside G-tilde. Used inside the counted proof at lines 516-529. Supplies the quotient and pull-back constructions that give the two inverse correspondences.
\end{itemize}\medskip
	\begin{Proposition}\label{extension}
		The $G$-equivariant Stolz exact sequence associated to $X$ is isomorphic to
		the $\widetilde{G}$-equivariant Stolz exact sequence associated to
		$\tilde{X}$.
	\end{Proposition}

	\begin{proof}
		The action of $\pi_1(X)\subset \widetilde{G}$ on $\tilde X$ is
		free. Consequently, for every $\widetilde{G}$-map ${f\colon \bar M\to\tilde X}$ the
		action of $\pi_1(X)$ on $\bar M$ is free. We can therefore quotient out this
		action and obtain a~cycle~${M:=\bar M/\pi_1(X)\to X}$ with residual action of
		$G:=\widetilde{G}/\pi_1(X)$.
		Vice versa, given a~$G$-map~${f\colon M\to X}$ we can pull back $\tilde X\to
		X$ along $f$ and obtain a $\pi_1(X)$-covering $\tilde M\to M$ with map
		$\tilde f\colon \tilde M\to \tilde X$. Because $f$ is a $G$-map, it is
		straightforward to pull back
		also the action of $\widetilde{G}$ to an action on $\tilde M$ which covers
		the action of $G$ on $M$ and such that $\tilde f$ is a $\widetilde{G}$-map
		(indeed, the $\widetilde{G}$-action on $\tilde M$ is defined mapping a point
		$(m,\tilde x)\in \tilde M\subset M\times \tilde X$ by
		$(\alpha,g)\in\widetilde{G}$ to $(gm,\alpha(\tilde x))$).
		
		The same construction gives a bijection between $G$-invariant metrics on $M$
		and $\widetilde{G}$-invariant metrics on $\tilde M$, and also works for
		bordisms.
		These constructions are clearly inverse to each other and preserve all
		additional structure and define the required bijections.
	\end{proof}

\end{document}
